% ============================================================
% v2/sections/related.tex  関連研究（短縮・再構成版 2026-10-27）
% ============================================================
\section{関連研究}\label{sec:related}

\paragraph{シミュレーション説と自他の共有表現}
神経科学では，自分の行為・感覚と他者のそれとに共通する表現と，両者を区別する仕組みの両方が社会的認知に必要だと整理されている\cite{decety2003shared}．
計算論の側では，運動制御のための内部モデルを他者の行為の理解に用いる枠組み\cite{wolpert2003unifying}や，ミラーニューロン系の働きを予測符号化として捉える考え方\cite{kilner2007predictive}が提案されてきた．
Gordon\cite{gordon2021simulation}は，自分の意思決定の仕組みが他者の行動の仮説を試すために使われ，そこでの信念や欲求は自他を区別しない符号化で表されると論じた．
重ね合わせ機構\cite{noguchi2022superposition}は，こうした考えを，自他の区別を与えない予測学習の計算モデルとして具体化したものと位置づけられる．
本研究は，この機構の上で，自分の価値に関わる予測が他者に使われるための条件を調べる．

\paragraph{投影と調整}
人は，自分の判断や好み，その時の身体の状態を他者に当てはめて過大に見積もる\cite{ross1977false,nickerson1999how,vanboven2003social}．
他者の視点を取ることは，自分の視点を出発点にして調整する過程として説明され\cite{epley2004perspective}，他者の好みを推し量るときの内側前頭前野の活動も，自分と他者の隔たりに応じて変わる\cite{tamir2010neural}．
Suzuki ら\cite{suzuki2012learning}は，他者の価値に基づく選択を予測するとき，自分の価値づけの信号が流用されること，他者の行動の予測誤差による補正が他者の価値が自分と異なるときに重要になることを示唆した．
本研究のモデルは，この区別でいえば自分の価値づけを流用する経路を持ち，他者の価値を補正する専用の経路を持たない．

\paragraph{自分をもとに他者を予測する計算モデル}
Raileanu ら\cite{raileanu2018modeling}は，自分の目標を入力の一つとする自分の方策を他者に当てはめ，他者の行動がもっともらしくなるように目標の入力を最適化して他者の目標を推した．
そこでは自分のモデルは報酬に基づく強化学習で得た方策で，目標が入力として明示されている．
本研究の見込みは，報酬も目標の入力も用いない予測学習の中で，自分の行動の経験から生まれる．
Copete ら\cite{copete2016motor}は，一つの予測器を自分の行為と他者の行為の観察の両方に用い，自分の運動の学習が他者の行為の結果の予測を助けることを示したが，観察者と行為者が同じ視点を共有すると仮定して自他の対応を与えている．
重ね合わせ機構では他者の内部状態が自分とは別に保たれるので，自他の重みの共有だけを外す対照や，学習後に自分側の重みを写す比較を作ることができる．
本研究はこれを用いて，他者の予測に使われるのが何かを切り分ける．
他者の行動から目標や好みを推す問題は，逆計画\cite{baker2009action,baker2017rational}や逆強化学習\cite{ng2000algorithms}として広く研究されている．
本研究の照合も候補ごとの予測を観察と比べる構造を持つが，推論の手法そのものを新しいとは主張しない．
